

Suppose that an agent's behaviour is governed by a Markov control model $(X,A,\{A(x) \mid x\in X\},Q,\elltrue)$ \citep[Def.~2.2.1]{savorgnan2009discrete}, where
\begin{itemize}
\item $X\subset\mR^{n_x}$, $A\subset\mR^{n_a}$ are the Borel space for state and action, respectively;
\item $\{A(x) \mid x\in X\}\subset A$ is the set of feasible control. Further, let $\mK:=\{(x,a) \mid x\in X,a\in A(x)\}\subset X\times A$ be the set of feasible state-action pairs;
\item $Q:(\mathcal{B}(X)\times \mK)\rightarrow [0,1]$ is the transition kernel;
\item $\elltrue:\mK\rightarrow\mR$ is the running cost.
\end{itemize}
Further, let $(\Omega,\mathcal{F},\mP)$ be the probability space that carries the state-action sequence $\{(\mfx_t,\mfa_t)\}_{t=1}^\infty$.
We assume that the agent make its decision by seeking to optimize the following infinite-horizon expected total discounted cost
\begin{align}
\min_{\pi}J(\pi;\mu_1)=\mE\left[\sum_{t=1}^\infty\alpha^{t-1}\elltrue(\mfx_t,\mfa_t)\right],\label{eq:forward_problem_obj_fun}
\end{align}
where $\alpha\in(0,1)$ is the decay constant, and the initial value $x_1$ is distributed according to $\mu_1$ on $X$.
In particular, the optimal policy $\bm{\pi}=\{\pi_t\}$ is a sequence of stochastic kernels $\pi_t:A\times (\mK^{t-1}\times X)\rightarrow [0,1]$, and the action $a_t$ that the agent taken at each time step $t$ is sampled from the distribution $\pi_t(a|x_1,a_1,\cdots,x_{t-1},a_{t-1},x_t)$.

To proceed, we make the following assumptions.
\begin{assumption}[Markov transition kernel]\label{ass:transition_kernel}
 The Markov transition kernel is a priori known. In addition, the Markov transition kernel $Q(\cdot|x,a)$ is Lipschitz continuous in total-variation distance, i.e., there exists $L_Q\ge 0$ such that for all $\eta_1:=(x_1,a_1)\in\mathbb{K},\eta_2:=(x_2,a_2)\in\mathbb{K}$, 
\begin{align*}
\|Q(\cdot|\eta_1)-Q(\cdot|\eta_2)\|_{\tv}\le L_Q\|\eta_1-\eta_2\|_2.
\end{align*}
\end{assumption}
\begin{assumption}[State-action set hyper-cube]\label{ass:compact_support}
The state-action set $X\times A$ is a hyper cube, namely, $X=\times_{i=1}^{n_x}[x_{i,\min},x_{i,\max}]$ and $A = \times_{i=1}^{n_a}[a_{i,\min},a_{i,\max}]$. Moreover, for all $x\in X$, it hold for the feasible action set $A(x)$ that $A(x)=A$. Consequently, $\mK = X\times A$.
\end{assumption}
\begin{assumption}[Cost function]\label{ass:obj_func}
  The cost function $\elltrue$ has prior structures in the sense that it belongs to a function class that can be parameterized linearly by $\thetatrue_\ell\in\mR^{n_\ell}$, i.e., $\elltrue(\cdot,\cdot) = \thetatrue_\ell^T\varphi(\cdot,\cdot)$, where $\varphi(\cdot,\cdot) = [\varphi_1(\cdot,\cdot),\ldots,\varphi_{n_\ell}(\cdot,\cdot)]^T$ is the vector of (known) $L_\varphi$-Lipschitz continuous feature functions, $\elltrue \in \lip(L_\ell)$ where $L_\ell = \|\thetatrue\|_1 L_\varphi$.
\end{assumption}
Moreover, under the Assumptions \ref{ass:transition_kernel}, \ref{ass:compact_support}, \ref{ass:obj_func}, we can further reformulate the problem of optimizing \eqref{eq:forward_problem_obj_fun} as a linear optimization problem in the infinite-dimensional space of measures \citep[Chap.~6]{hernandez2012discrete}.
In particular, given a policy $\bm{\pi}$ and the initial value's distribution $\mu_1$, the state-control occupation measure at time step $t$ is defined as
\begin{align*}
  \mu_t^\pi(\Gamma_x\times \Gamma_a) :=\mE[\indicator_{\Gamma_x\times \Gamma_a}(\mfx_t,\mfa_t)],\quad \forall (\Gamma_x\times\Gamma_a)\subset\mK.
\end{align*}
In addition, let $\proj:\mathcal{M}_+(\mK)\rightarrow\mathcal{M}_+(X)$ be the linear projection operator that transform any probability measure $\mu\in\mathcal{M}_+(\mK)$ to its marginal distribution $\tilde{\mu}\in\mathcal{M}_+(X)$ on the state. 
Moreover, for all $\Gamma_x\subset X$, let $q:\mathcal{M}_+(\mK)\rightarrow\mathcal{M}_+(X)$ be the linear forward time shift operator defined by
\begin{align*}
  (q\mu_t^\pi)(\Gamma_x) = \int_{\mK} Q(\Gamma_x|x,a)\mu_t^\pi(\mathrm d x,\mathrm d a).
\end{align*}
Therefore, in view of the definition of the transition kernel $Q$ and the state-control occupation measure, it is clear that $\proj\mu_{t+1}^\pi = q\mu_t^\pi$. Let us further define the following occupation measure
\begin{align}\label{eq:mu^pi_def}
  \mu^\pi:=\sum_{t=1}^\infty \alpha^{t-1} \mu_t^\pi.
\end{align}
Equipped with the set-up above, we can reformulate the the problem of optimizing \eqref{eq:forward_problem_obj_fun} as
\begin{subequations}\label{eq:forward_problem_measure_form}
\begin{align}
  \min_{\mu^\pi\in\mathcal{M}_+(\mK)} &\;\;\langle \elltrue,\mu^\pi\rangle \label{eq:forward_problem_measure_form_obj}\\
  \st &\;\;(\proj-\alpha q)\mu^\pi = \mu_1.\label{eq:forward_problem_measure_form_dynamics}
\end{align}
\end{subequations}
In fact, the Assumptions \ref{ass:transition_kernel}, \ref{ass:compact_support}, \ref{ass:obj_func} guarantees that \eqref{eq:forward_problem_measure_form} is equivalent to optimizing \eqref{eq:forward_problem_obj_fun} \citep[Thm.~6.3.7]{hernandez2012discrete} in the sense that there is an equivalence between the feasible points of the two problems and that they both attains the same optimal value.
\footnote{More specifically, Assumption \ref{ass:transition_kernel}, \ref{ass:compact_support} and \ref{ass:obj_func} together imply \citep[Ass.~4.2.1]{hernandez2012discrete}.}
The equivalence between \eqref{eq:forward_problem_obj_fun} and \eqref{eq:forward_problem_measure_form} means that we can consider the latter formulation when we consider the IOC problem to come. 

To this end, the optimality conditions for the above problem are given in \eqref{eq:opt_cond_inf_dim}, and they can derived as follows. We first relax the constraint \eqref{eq:forward_problem_measure_form_dynamics} with a multiplier $\Vtrue \in \mathscr{F}(X)$, and get the Lagrangian $\mathcal{L} : \mathcal{M}_+(\mK) \times \mathscr{F}(X) \rightarrow \mR$
\begin{align*}
\mathcal{L}(\mu^\pi, \Vtrue) & = \langle \elltrue,\mu^\pi\rangle - \langle \Vtrue, (\proj-\alpha q)\mu^\pi -\mu_1 \rangle \\
& =  \langle \elltrue + (\alpha q^*-\proj^*)\Vtrue,\mu^\pi\rangle + \langle \Vtrue, \mu_1 \rangle.
\end{align*}
Considering the first term in the above equation, we note that $\inf_{\mu^\pi\in\mathcal{M}_+(\mK)}  \mathcal{L}(\mu^\pi, \Vtrue) > -\infty$ only if 
\begin{subequations}\label{eq:opt_cond_inf_dim}
\begin{equation}
\elltrue+(\alpha q^*-\proj^*)\Vtrue\ge 0. \label{eq:bellman_inequality}
\end{equation}
On the other hand, since the above equation holds, the infimum of the Lagrangian is attained for any $\mu^\pi$ such that
\begin{equation}\label{eq:complementary_slackness}
  \langle \elltrue + (\alpha q^*-\proj^*)\Vtrue,\mu^\pi\rangle = 0,
\end{equation}
for example $\mu^\pi = 0$. However, the latter is not primaly feasible, which gives us the last optimality condition, namely
\begin{align}
  \eqref{eq:forward_problem_measure_form_dynamics},\quad \mu^\pi\in\mathcal{M}_+(\mK).\label{eq:primal_feasible}
\end{align}
\end{subequations}
In these equations, the adjoint operators are defined by $q^*\Vtrue:\mK\mapsto\mR$ and $\proj^* \Vtrue:\mK\mapsto\mR$ via
\begin{subequations}  \label{eq:adjoints}
  \begin{align}
    (q^*\Vtrue)(x,a) &=\int_X \Vtrue(\chi)Q(d\chi|x,a)\label{eq:cost_to_go}\\
    (\proj^*\Vtrue)(\mfx,\mfu) &= \Vtrue(\mfx)\indicator_{\mK}(\mfx,\mfa)=\Vtrue(\mfx).\label{eq:value_func}
  \end{align}
\end{subequations}
We further assume that as an observer, we can observe $M$ trials of the agent.  More precisely, let the trajectories $(\mfx_t^i, \mfa_t^i)_{t=1}^{N}$  be Independently Identically Distributed (I.I.D) over $i=1:M$. The observed trials $\{(x_t^i,a_t^i)\}_{i=1}^M$ are the realizations of the I.I.D state-action pairs $\{(\mfx_t^i,\mfa_t^i)\}_{i=1}^M$.
Moreover, we make the following assumptions on the initial value distribution $\mu_1$ and the observed agent's optimal actions.

\begin{assumption}[Initial value distribution]\label{ass:init_value}
The support of the initial value distribution $\mu_1$ is $X$.
\end{assumption}

\begin{assumption}[Deterministic stationary policy]\label{ass:stationary_markov_policy}
The observed agent's optimal actions $\{a_t^i\}_{i=1}^M$ for all $t=1:N$ are generated from a deterministic stationary policy \citep[Def.~2.3.2]{hernandez2012discrete}. Namely, it holds for the optimal policy $\bm{\pitrue}=\{\pitrue_t\}$ that $\pitrue_t(a|x_1,a_1,\cdots,x_{t-1},a_{t-1},x_t) = \delta(\pitrue(x_t))$, which is the Dirac measure concentrated at the measurable function $\pitrue:X\rightarrow A$.
\end{assumption}

With the above set-up, we are ready to present the problem that is considered in this paper. 

\begin{problem}\label{pro:ioc}
  Given the $M$ observations of the optimal ``state-action pair" trajectories $\{(x_t^i,a_t^i)\}_{t=1}^N$ that are generated by \eqref{eq:forward_problem_obj_fun}, construct a consistent estimator of $\pitrue:X\rightarrow A$ that takes the structure of $\elltrue$ into account.
\end{problem}



Note that different cost functions may lead to identical optimal policies.
For instance, $\pi(x)=0$ is optimal to both $\ell_1(x,a)=a^2$ and $\ell_2(x,a)=a^4$.
Hence, $\thetatrue_\ell$ is in general not identifiable unless the feature functions vector $\varphi$ is carefully chosen.
Nevertheless, motivated by the practical need of ``learning from demonstration", in this paper we choose to not focus on the identifiability of IOC. Instead, the question of identifiability of the cost function is left for future research.
Despite this, we can still do consistency analysis in the sense that the collected state-action pairs $(x_t^i, a_t^i)$ are indeed optimal to the reconstructed cost function as the data amount tends to infinity. 
Moreover, we can also get an estimate $\hat{\pi}$ of the underlying optimal policy $\pitrue$ that corresponds to the estimated cost $\hat{\theta}_\ell^T\varphi$ (cf.~\citep{zhang2024inverse}).
Such policy estimation intrinsically takes the prior knowledge of the structures in $\elltrue$ into account, in contrast to ``behaviour cloning" \citep{sammut2017behavioral, codevilla2019exploring, torabi2018behavioralcloningobservation}. More precisely, behaviour cloning parameterizes a family of deterministic stationary policies $\pi(\cdot;\theta)$ (typically a neural network) and solves the following optimization problem
\begin{align}
\min_{\theta}\frac{1}{M}\sum_{i=1}^M\sum_{t=1}^N\|\mfa_t^i-\pi(\mfx_t^i;\theta)\|_2^2.
\label{eq:behaviour_cloning}
\end{align}
Indeed, it also offers a consistent estimator of $\pitrue$, yet such control policy reconstruction is vulnerable to out-of-distribution data samples.
Such samples can occur, e.g., due to compounding errors when using the control policy \citep{ross2011reduction}, due to distribution shifts in the noise \citep{pmlr-v70-pinto17a}, or when trying to generalize the behaviour to new environments \citep{pmlr-v48-finn16, fu2017learning}. 
In contrast, by taking the structures in $\elltrue$ into account, we obtain a more robust control policy estimate.
